\subsection*{What the architecture contributes, and what its attention does not explain}

Two questions remain about the model that produced the findings above: whether its neural half
contributes anything beyond the classical arm, and whether its attention and gating weights explain
its predictions. The first has a small positive answer, the second a negative one.

\textbf{The neural branch adds a small, seed-stable increment.} On the lineage holdout the full
stack exceeded the stacked classical arm by $+0.028$ [0.006, 0.052] mean \nap{}, clustered by
sequence type. Because that interval is narrow and comes from one network initialisation, we
retrained the neural branch under two further seeds, holding the fold partition, gene set,
hyperparameters and classical arm fixed so that only weight initialisation, dropout and batch order
changed. The margin was $+0.020$ [0.003, 0.042] and $+0.018$ [$-0.001$, 0.037] (mean $+0.022$, SD
0.005), which passes the pre-registered rule (positive mean, interval above zero in two of three
seeds). The neural branch alone stayed well below the classical arm in every seed (0.405--0.435
against 0.556). The architecture is one small reproducible contribution to a stack, not the source
of this paper's findings.

\textbf{The gate is readable but does not route by mechanism.} The gate's design rests on a
prediction: a drug whose resistance comes from an acquired gene should route to the
presence/absence branch, and a drug whose resistance comes from a point mutation in a target enzyme
to the motif branch. The presence/absence branch does dominate, taking 0.553 to 0.807 of the weight
for every drug, which is consistent with how much \Ecoli{} resistance is carried on acquired genes
(Table~\ref{tab:gate}). But ciprofloxacin is the panel's clearest point-mutation drug, the one
whose known targets our own attribution analysis recovers, and it gives the motif branch 0.081 ---
the lowest weight in the table. The prediction fails on the drug where it was most testable. The
cross-antibiotic branch takes its highest weight there (0.367), which cannot be read as mechanism
either: ciprofloxacin is the only fluoroquinolone in the panel and has no same-class partner to
share with, and its co-resistance link to ceftazidime ($\phi = 0.52$ in this collection) is
co-occurrence. The gate tells us which branch the model leans on, which is worth knowing, and not
why.

\begin{table}[t]\centering\small
\caption{\textbf{Mean gate weights on the random-split test set.} Each row sums to one. Branch A is
the motif encoder, B the gene attention module, C the cross-antibiotic module.}
\label{tab:gate}
\begin{tabular}{lrrr}
\toprule
Drug & Branch A (motif) & Branch B (presence/absence) & Branch C (cross-antibiotic) \\
\midrule
AMP & 0.153 & 0.649 & 0.198 \\
CIP & 0.081 & 0.553 & 0.367 \\
AMC & 0.098 & 0.641 & 0.261 \\
SXT & 0.100 & 0.807 & 0.093 \\
GEN & 0.129 & 0.676 & 0.195 \\
CTX & 0.076 & 0.735 & 0.189 \\
CAZ & 0.068 & 0.748 & 0.183 \\
\bottomrule
\end{tabular}
\end{table}

\textbf{Attention recovers no known resistance gene; post-hoc attribution on the classical arm
recovers eight.} We collected the most influential genes from two sources, the attention weights of
the three branches and SHAP values \cite{lundberg2017shap} on the five classical learners, giving
1,143 marker entries, and matched their Prokka \cite{seemann2014prokka} locus tags against the
eggNOG \cite{huertacepas2019eggnog} functional annotation shipped with the collection. We annotated
836 (73.1\%), reporting the rest as unannotated rather than guessing. Eight annotated markers are
known resistance genes: \textit{gyrA} and \textit{parC}, the fluoroquinolone targets;
\textit{emrB}, \textit{mdtC}, \textit{mdtD} and \textit{mdtO}, components of multidrug efflux
pumps; \textit{cusA}, which belongs to the copper and silver efflux system and is therefore a
heavy-metal rather than a drug transporter, so we count it as efflux and not as evidence of an
antibiotic mechanism; and one locus annotated as a $\beta$-lactamase.

Recovering \textit{gyrA} and \textit{parC} unprompted is the one gene-level result that the lineage
holdout strengthens rather than weakens. These are the textbook targets of ciprofloxacin
resistance, and ciprofloxacin is the drug that survived the holdout best, losing 0.026 once
prevalence is matched against up to 0.239 for the others. A signal that is a small, well-defined
mutation is exactly what a model can learn once and apply in an unrelated isolate.

All eight came from SHAP on the classical learners. None came from the attention weights of the
three branches, although 420 of the 1,143 marker entries were selected by them; the closest the
attention came was \textit{cusC}, the outer-membrane partner of the same heavy-metal efflux system,
surfacing for three drugs in the cross-antibiotic branch. Two explanations besides a genuine failure
could produce that zero, and both were tested. Annotation coverage does not account for it:
coverage was 90.7\% for the motif branch and 70.0\% for the cross-antibiotic branch, and neither
produced a resistance gene, while the presence/absence branch sits lower at 60.0\% --- which is its
own problem, since the branch the gate relies on most is the one we can interpret least. Nor is it
an artefact of list depth: repeating the annotation on the top 100 genes per branch per drug rather
than the top 20, a fivefold widening, raised coverage but recovered no resistance gene.

The attention weights are also diffuse. No single locus in the presence/absence branch exceeds
0.015 of the attention mass for any drug, and in the other two branches none exceeds 0.005, so the
pooling layer behaves close to a mean over the genes a strain carries. What focus there is points at
provenance rather than function. Each pan-genome cluster carries the locus tag of the assembly it
was called from, which gives a cheap test of whether a branch's attended genes are drawn evenly
across the catalogue. For the presence/absence branch they are not: its top genes concentrate on
one source assembly at 3.7 to 6.7 times that assembly's share of the catalogue for six of the seven
drugs (significant at 5\% for four of them), and at 0.93 times for the seventh. These counts rest
on fewer than a dozen tagged genes per drug, so we read them as the direction of the lens rather
than a formal test.

\textbf{Strain conditioning helps where it is needed and costs interpretability.} The degenerate
scoring function described in Methods is a concrete defect with a concrete fix, so we ablated the
fix, training in both directions with identical seeds, folds, splits and hyperparameters and
changing only that term (Table~\ref{tab:attnfix}). It buys nothing where it is easiest to measure:
$-0.002$ on the random split and $+0.006$ on the full stack under the lineage holdout, both well
inside the seed variability established above. It buys $+0.009$ on the neural branch alone under
that holdout, where the defect does its damage and where the classical arm does not move at all.
And three of the four interpretability measures move against it: the resistance information of the
attended genes falls from 0.047 to 0.021, their lineage determination rises from 0.277 to 0.430,
and the single resistance gene the standard formulation had recovered is lost. The one measure that
improves is the provenance signature, which weakens from 20.5-fold to 4.7-fold enrichment on one
assembly.

\begin{table}[t]\centering\small
\caption{\textbf{Ablation of strain conditioning in the gene attention module.} Everything but that
one term is identical; ``conditioned'' is the delivered model reported everywhere else. $U_{ST}$ is
the fraction of a gene's presence/absence entropy explained by sequence type and $U_{res}$ the
fraction explained by resistance, over the top twenty genes per drug; lower $U_{ST}$ and higher
$U_{res}$ are better.}
\label{tab:attnfix}
\begin{tabular}{lrrrl}
\toprule
Measure & Standard & Conditioned & Change & Reference \\
\midrule
\multicolumn{5}{l}{\textit{Predictive performance (mean AP)}} \\
Random split, full stack & 0.853 & 0.850 & $-0.002$ & --- \\
Phylogroup-stratified split & 0.857 & 0.856 & $-0.001$ & --- \\
Lineage holdout, full stack & 0.640 & 0.645 & $+0.006$ & --- \\
Lineage holdout, classical arm alone & 0.622 & 0.622 & $\pm 0$ & --- \\
Lineage holdout, neural branch alone & 0.511 & 0.520 & $+0.009$ & --- \\
\midrule
\multicolumn{5}{l}{\textit{Clinical safety and interpretability}} \\
Mean very major error rate & 0.239 & 0.241 & $+0.002$ & --- \\
Resistance genes recovered by attention & 1 & 0 & $-1$ & SHAP: 8 \\
$U_{res}$ of attended genes & 0.047 & 0.021 & $-0.025$ & SHAP: 0.048 \\
$U_{ST}$ of attended genes & 0.277 & 0.430 & $+0.153$ & random gene: 0.539 \\
Enrichment on one source assembly & $20.5\times$ & $4.7\times$ & $-15.8\times$ & --- \\
\bottomrule
\end{tabular}
\end{table}

We report this as a trade rather than an improvement, and as a correction to an earlier run of the
same ablation in which the term appeared to be worth $+0.035$ on the neural branch and to improve
three interpretability measures. Repeating the protocol reversed most of those signs. The lesson
is the one the seed experiment above makes for the stack's own margin: an effect of this size is not
separable from training variability without repetition, and we do not claim it as one.

\textbf{Attribution does not agree with itself.} Two statistics bound how far any of the gene-level
results can be pushed. The attended gene sets are not lineage markers --- their $U_{ST}$ of 0.430
is below a random gene's 0.539, and no draw among 500 random sets of equal size fell below them ---
but they carry little resistance information either, at $U_{res} = 0.021$ against 0.048 for the
SHAP-selected genes, a factor of 2.25. The presence/absence lens did not read the wrong thing so
much as read almost nothing, and the classical learners that generalise better also pick the more
informative genes. Separately, the five classical learners disagree among themselves: across all 70
(drug, learner-pair) combinations the mean Spearman correlation of SHAP rankings is 0.061 and the
mean top-twenty Jaccard overlap 0.180, so two attributions of the same trained pipeline on the same
features share three or four genes in twenty. Among the full set of 32,441 loci, a random gene's
presence is explained to 54\% by sequence type, which is the single clearest reason a
lineage-sensitive split was necessary in the first place. We therefore report the gene-level
findings as leads, informative where they agree with independently established biology, and not as
evidence that any attribution method has identified what the model uses.
